\documentclass{article}

% "preprint" option is used for arXiv or other preprint submissions
\usepackage[preprint]{neurips_2025}

\usepackage[utf8]{inputenc}
\usepackage[T1]{fontenc}
\usepackage[spanish]{babel}
\usepackage{hyperref}
\usepackage{url}
\usepackage{booktabs}
\usepackage{amsfonts}
\usepackage{amsmath}
\usepackage{nicefrac}
\usepackage{microtype}
\usepackage{xcolor}
\usepackage{array}
\usepackage{multirow}

\title{No Free Lunch: Hedge Fund Time Series Forecasting}

\author{%
  Andrés Felipe Mosquera Hernández \\
  Universidad de los Andes \\
  \texttt{a.mosquerah2@uniandes.edu.co} \\
  \And
  Josue Briceño Urquijo \\
  Universidad de los Andes \\
  \texttt{j.bricenou@uniandes.edu.co} \\
  \And
  Michael Javier Patiño Pantoja \\
  Universidad de los Andes \\
  \texttt{mj.patino@uniandes.edu.co} \\
}

\begin{document}

\maketitle

% \begin{abstract}
%   TODO: escribir el abstract una vez finalizado el trabajo.
% \end{abstract}


% ============================================================
% 1. INTRODUCCIÓN
% ============================================================

\section{Introducción}

El pronóstico de series de tiempo es una tarea fundamental en las finanzas cuantitativas, siendo base de la asignación de portafolios, la gestión de riesgo y las estrategias de trading sistemático. Las competencias M4 \citep{makridakis2020m4} y M5 \citep{makridakis2022m5} han demostrado que los métodos modernos especialmente los árboles de gradient boosting \citep{ke2017lightgbm} y las arquitecturas híbridas superan sistemáticamente las líneas base estadísticas clásicas cuando se dispone de datos etiquetados a gran escala. La experiencia de soluciones ganadoras en Kaggle confirma además que las estrategias de validación cruzada temporal y los ensambles de predicciones recursivas y directas son decisivos para lograr una generalización robusta fuera de muestra \citep{m5first2020, kaggleguide}.

En la práctica, sin embargo, las series financieras presentan dificultades adicionales: un ratio señal-ruido extremadamente bajo, no estacionariedad y dinámicas heterogéneas. La competencia \emph{Hedge Fund Time Series Forecasting} \citep{adatacompany2026} exige abordar este escenario mediante modelos predictivos sobre datos jerárquicos anonimizados, prohibiendo además el uso de información externa. Debido a una evaluación basada en predicción secuencial estricta, el uso de datos futuros (fuga de datos) es severamente penalizado. El desafío principal consiste en generalizar adecuadamente sin sobreajustarse al leaderboard público, el cual resulta ruidoso.

% -----------------------------------------------------------
\subsection{Justificación}

El presente trabajo se justifica en tres dimensiones fundamentales. En términos de \textbf{relevancia práctica}, las mejoras marginales en la precisión del pronóstico se traducen en retornos ajustados por riesgo fuertemente superiores de cara a fondos de cobertura reales ($>\$4$ billones USD bajo gestión); por lo tanto, trabajar con estas severas restricciones produce una metodología directamente transferible a entornos de producción bajo causalidad estricta. Respecto al \textbf{vacío de conocimiento}, la mayoría de la investigación actual requiere características interpretables (ej. OHLCV, macroeconomía), mientras que aquí evaluaremos si un aprendizaje puramente guiado por los datos puede recuperar una estructura predictiva sin conocimiento explícito de dominio. Finalmente, sobre la \textbf{pertinencia metodológica}, el alto riesgo de sobreajuste al leaderboard público exige estructurar un escrupuloso framework de validación cruzada temporal reproducible y escalable.

% -----------------------------------------------------------
\subsection{Objetivos}

\textbf{Objetivo general.} Desarrollar y evaluar una metodología de machine learning para el pronóstico de series de tiempo financieras anonimizadas bajo restricción estricta de predicción secuencial, maximizando el \emph{weighted RMSE skill score} en el test privado de la competencia Kaggle \emph{Hedge Fund -- Time Series Forecasting}.

Para cumplir con el propósito principal, se establecen los siguientes \textbf{objetivos específicos}:

\begin{itemize}
    \item \textbf{Construcción causal:} Diseñar un pipeline de ingeniería de features que calcule todos los estadísticos rolling y de rezago exclusivamente sobre el pasado ($\texttt{ts\_index} \leq t$), con el criterio de éxito de probar matemáticamente y computacionalmente la ausencia de fuga de datos.
    \item \textbf{Explotación jerárquica:} Analizar la estructura de los identificadores (\texttt{code}, \texttt{sub\_code}, \texttt{sub\_category}, \texttt{horizon}) para identificar patrones y dependencias transversales, apuntando a superar a las líneas base simples que analizan cada grupo y horizonte de manera aislada.
    \item \textbf{Entrenamiento y validación:} Comparar distintos enfoques de modelado —regresión regularizada, gradient boosting y modelos de secuencias— bajo un esquema de validación cruzada temporal estricta, seleccionando el modelo con mayor estabilidad entre folds locales.
    \item \textbf{Control de sobreajuste:} Evaluar y descartar el ruido estandarizado del \emph{leaderboard} público midiendo el impacto del \textit{overfitting}, para lograr efectuar la final selección de modelos basados rigurosamente en la estabilidad, en particular, priorizando reducciones en la varianza o desviación estándar entre \emph{folds} locales.
\end{itemize}


% ============================================================
% 2. ANTECEDENTES Y TRABAJOS RELACIONADOS
% ============================================================

\section{Antecedentes y Trabajos Relacionados}

La literatura relevante para esta competencia cubre cinco áreas: (i)~arquitecturas de deep learning para pronóstico de series de tiempo (TSF), (ii)~gradient boosting y modelos globales, (iii)~competencias de pronóstico, (iv)~pronóstico financiero, y (v)~generalización bajo cambio de distribución.

% -----------------------------------------------------------
\subsection{Tres trabajos relevantes}

\paragraph{TimeXer: Empowering Transformers for Time Series Forecasting with Exogenous Variables \citep{wang2024timexer}.}
El paper propone un modelo basado en la arquitectura Transformers pero adaptada específicamente para mejorar la calidad del forecast en series temporales en los que intervienen variables exógenas que influyen en la variable objetivo. El problema que aborda el paper es que muchos modelos actuales son capaces de manejar bien la dependencia temporal en una serie o entre varias variables, pero rara vez manejan ambas propiedades simultáneamente. Esto limita su desempeño en escenarios reales en los que los factores externos influyen en la evolución de la serie temporal. Además, muchas variables exógenas presentan irregularidades en los datos, tales como valores faltantes o diferentes frecuencias, lo que dificulta la integración en los modelos de predicción.

Para superar estos problemas, los autores proponen, TimeXer, que es un modelo que aprovecha la arquitectura tradicional del modelo Transformers pero con la estrategia de representación jerárquica para manejar variables endógenas y exógenas de manera diferente. La serie objetivo se divide en \textit{patches} temporales para procesar la dependencia temporal interna. Las variables exógenas, por su parte, son representadas como tokens en el nivel de variable, facilitando el tratamiento de datos irregulares o de longitud variable. Finalmente, un token global aprendible conecta las dos representaciones, facilitando el influir de la información exógena en las secciones temporales pertinentes a través de mecanismos de \textit{cross-attention}. Con esta estructura, el modelo logra capturar tanto las dependencias temporales dentro de la serie como las relaciones entre las variables exógenas y la serie objetivo.

Los experimentos se han llevado a cabo en varios datasets reales de forecasting, incluyendo datos meteorológicos, de tráfico, eléctricos, y de precios eléctricos. Los experimentos muestran claramente que el modelo supera a varios modelos de referencia, tanto en escenarios univariados como en escenarios multivariados, incluyendo a modelos Transformer, así como a modelos basados en RNN, grafos, etc. Además, el modelo ha mostrado su capacidad para tratar adecuadamente la información exógena, para escalar a datasets de gran magnitud, así como para tratar escenarios complejos como ausencias de datos, o variables de diferentes frecuencias. En resumen, el estudio concluye que el tratamiento adecuado de las variables exógenas a través de esta estructura mejora significativamente la precisión en el forecasting en escenarios reales de series temporales.

\paragraph{Are Self-Attentions Effective for Time Series Forecasting? \citep{kim2024cats}.}
Este paper presenta CATS (Cross-Attention-only Time Series transformer), una arquitectura de Transformers al prescindir totalmente de las capas de auto-atención. Esta paper presenta que la auto-atención es intrínsecamente ineficiente para el pronóstico de series temporales debido a su invariancia de permutación, lo que dificulta la captura de dependencias secuenciales. En su lugar, CATS utiliza exclusivamente un mecanismo de atención cruzada, donde la serie de tiempo de entrada se codifica como keys y values, mientras que los queries se definen como parámetros aprendibles dependientes del horizonte futuro.

Técnicamente, el modelo demuestra una capacidad superior para capturar tanto la periodicidad de la señal como las perturbaciones (shocks) o cambios bruscos, superando las limitaciones de la auto-atención en el modelado de estos patrones críticos. Además, la arquitectura optimiza el consumo de memoria y la carga computacional mediante una estrategia de compartición de parámetros entre diferentes horizontes y el uso de técnicas de patching. En términos de métricas de rendimiento como el error cuadrático medio (MSE) y el error absoluto medio (MAE), CATS logra establecer un nuevo estado del arte, superando a modelos prominentes como TimeMixer y PatchTST en horizontes de predicción que van desde los 96 hasta los 720 pasos de tiempo.

\paragraph{LightGBM \citep{ke2017lightgbm}.}
Se propone LightGBM con dos innovaciones algorítmicas independientes. GOSS (\textit{Gradient-based One-Side Sampling}) retiene todas las instancias con gradiente elevado, las menos ajustadas por el modelo y submuestrea aleatoriamente las de gradiente bajo, aplicando un factor de amplificación para compensar el sesgo introducido; de este modo, el cómputo de las ganancias de split se aproxima con alta fidelidad usando una fracción de los datos. EFB (\textit{Exclusive Feature Bundling}) identifica grupos de features dispersas que rara vez toman valores no-nulos simultáneamente y las agrupa en un único bundle mediante coloración greedy de grafos, reduciendo la dimensionalidad efectiva sin pérdida de información. En la competencia M5, las soluciones ganadoras basadas en LightGBM con features de rezago sistemáticas superaron a todos los enfoques estadísticos clásicos \citep{makridakis2022m5, m5first2020}.

La limitación estructural del modelo es que trata cada fila como una observación i.i.d.: no incorpora ningún mecanismo de orden temporal ni de causalidad. Toda la información secuencial debe inyectarse externamente como features (rezagos, medias móviles, cuantiles), transfiriendo la responsabilidad de la restricción causal íntegramente al pipeline de preprocesamiento.

En este trabajo, LightGBM actúa como motor de predicción global, pero el diferenciador crítico es la etapa previa de ingeniería de features causales, que construye todos los estadísticos rolling y de rezago bajo la restricción $\texttt{ts\_index} \leq t$. Esta separación de responsabilidades, causalidad en el pipeline, capacidad predictiva en el modelo, es la contribución central de esta metodología.

% ============================================================
% 3. PLANTEAMIENTO DEL PROBLEMA
% ============================================================

\section{Planteamiento del Problema. Formalización}

\textbf{Notación.} Asumiendo que $g = (c,s,k,h)$ identifica un grupo-horizonte y $t$ un índice temporal válido, denotemos $\mathbf{x}_{g,t} \in \mathbb{R}^{86}$ al vector de atributos, e $y_{g,t} \in \mathbb{R}$ a la variable objetivo correspondiente.

La \textbf{Restricción de Predicción Secuencial} (estricta) impone que todo modelo $f_\theta$ genere una predicción en base \emph{exclusiva} a eventos ocurridos hasta el índice temporal evaluado. Matemáticamente:
\[
  \hat{y}_{g,t} = f_\theta\!\left(\mathcal{D}_{\leq t}\right), \quad \text{donde} \quad
  \mathcal{D}_{\leq t} = \bigl\{(\mathbf{x}_{g',t'},\ldots) \mid t' \leq t,\;\forall g'\bigr\}
\]
donde $\mathcal{D}_{\leq t}$ incluye cálculos como promedios móviles, normalización o cuantiles referenciados a índices $\leq t$. Ignorar esto es un motivo directo de invalidación.

El \textbf{Problema Causal de Pronóstico} entonces consiste en entrenar el modelo $f_\theta$ respetando la causalidad anterior, con la meta de maximizar la siguiente métrica objetivo o \textit{Weighted RMSE Skill Score}:
\begin{equation}
  \theta^* = \arg\max_\theta \sqrt{1 - \mathrm{clip}\!\left(
    \frac{\sum_{i \in \mathcal{I}} w_i(y_i - \hat{y}_i)^2}
         {\sum_{i \in \mathcal{I}} w_i\,y_i^2},\;0,\;1\right)}
  \label{eq:score}
\end{equation}
evaluado estrictamente fuera de muestra sobre el conjunto privado ($\mathcal{I}_{privado}$), y donde $w_i > 0$ es el peso estipulado y $\mathrm{clip}(r,0,1)$ asegura puntuaciones en un rango acotado $[0, 1]$.

% ============================================================
% 4. DATOS Y TÉCNICA DE EXTRACCIÓN
% ============================================================

\section{Datos empleados}

Los datos provienen de la competencia Kaggle \emph{Hedge Fund Time Series Forecasting} \citep{adatacompany2026}, aportados por un fondo de cobertura activo. Se distribuyen en dos archivos Parquet: \texttt{train.parquet} (94 columnas, incluye target y pesos) y \texttt{test.parquet} (92 columnas, sin target ni pesos). Las observaciones de test corresponden a un periodo temporal estrictamente posterior al de entrenamiento, aunque las fechas exactas permanecen ocultas.

\begin{table}[h!]
  \caption{Diccionario de datos. El campo \texttt{id} concatena las cinco columnas clave con separador \texttt{\_\_}.}
  \label{tab:datos}
  \centering
  \begin{tabular}{llp{4.8cm}p{3.0cm}}
    \toprule
    Columna & Tipo & Descripción & Implicación para el modelo \\
    \midrule
    \texttt{id}            & str   & Concatenación: \texttt{code}, \texttt{sub\_code}, \texttt{sub\_category}, \texttt{horizon}, \texttt{ts\_index} & Clave única de submisión. \\
    \texttt{code}          & str   & Entidad primaria (anonimizado) & \multirow{3}{3cm}{Jerarquía inter-grupo explotable.} \\
    \texttt{sub\_code}     & str   & Subgrupo dentro de \texttt{code} & \\
    \texttt{sub\_category} & str   & Categoría amplia (anonimizada) & \\
    \texttt{ts\_index}     & int   & Eje temporal entero; sin calendario explícito & Orden causal estricto. \\
    \texttt{horizon}       & int   & Régimen: 1 (corto), 3 (medio), 10 (largo), 25 (extra-largo) & Etiqueta, \emph{no} desplazamiento. \\
    \texttt{weight}        & float & Peso por muestra para la función de pérdida & Solo \texttt{sample\_weight}; \textbf{nunca} feature. \\
    \texttt{y\_target}     & float & Variable objetivo continua (solo en train) & SNR bajo (${\approx}0.063$ baseline). \\
    \texttt{feature\_a} - \texttt{feature\_ch} & float & 86 features numéricas anónimas; admiten nulos & Imputación causal (\emph{forward-fill}). \\
    \bottomrule
  \end{tabular}
\end{table}


% ============================================================
% 5. METODOLOGÍA
% ============================================================

\section{Metodología}
\label{sec:metodologia}

La metodología sigue un pipeline modular de cuatro etapas, cada una alineada con los objetivos específicos: ingeniería de features causales, análisis de estructura jerárquica, comparación de enfoques de modelado, y control de sobreajuste.

% -----------------------------------------------------------
\subsection{Ingeniería de features causales}

Todos los estadísticos derivados —medias móviles, desviaciones estándar rolling, cuantiles y rezagos— se computan condicionados a $\texttt{ts\_index} \leq t$, garantizando la ausencia de fuga de datos. Los valores faltantes en las 86 features anónimas se imputan mediante \emph{forward-fill} causal, propagando únicamente observaciones pasadas. La validación de causalidad se realizará mediante pruebas unitarias que verifiquen que ningún cálculo referencia índices $> t$.

% -----------------------------------------------------------
\subsection{Análisis de estructura jerárquica}

Se explotará la jerarquía definida por los identificadores \texttt{code}, \texttt{sub\_code}, \texttt{sub\_category} y \texttt{horizon} para construir features de agregación entre grupos. El análisis exploratorio identificará dependencias transversales entre series de un mismo \texttt{code} o \texttt{sub\_category}, con el fin de capturar señales que los modelos que tratan cada serie de manera aislada no pueden aprovechar.

% -----------------------------------------------------------
\subsection{Enfoques de modelado}

Se compararán tres arquitecturas de predicción seleccionadas por su idoneidad en series de tiempo financieras:

\begin{itemize}
    \item \textbf{Regresión regularizada (Ridge/Lasso):} Línea base interpretable que cuantifica la linealidad residual en las features causales.
    \item \textbf{Gradient boosting (LightGBM \citep{ke2017lightgbm}):} Motor de predicción global con soporte nativo para pesos de muestra (\texttt{sample\_weight}), necesarios para el cómputo del \emph{Weighted RMSE Skill Score}.
    \item \textbf{Modelos de secuencias (CATS \citep{kim2024cats}):} Arquitecturas con dependencia temporal explícita, sin el sesgo de invariancia de permutación propio de la auto-atención estándar.
\end{itemize}

Los tres enfoques se entrenarán bajo la misma estrategia de validación cruzada temporal para garantizar comparaciones equitativas.

% -----------------------------------------------------------
\subsection{Métricas de evaluación}

Las siguientes métricas guiarán la selección y el ajuste de modelos a lo largo del proyecto:

\begin{itemize}
    \item \textbf{Weighted RMSE Skill Score} (Ecuación~\ref{eq:score}): métrica primaria de la competencia; se optimiza sobre el conjunto de prueba privado y se estima localmente en cada fold de validación.
    \item \textbf{RMSE por horizonte:} Desagregación del error cuadrático medio para cada horizonte $h \in \{1, 3, 10, 25\}$, permitiendo identificar los regímenes con mayor dificultad predictiva.
    \item \textbf{Error Absoluto Medio (MAE):} Complemento robusto frente a valores atípicos, útil dado el bajo SNR de las series financieras.
    \item \textbf{Desviación estándar del score entre folds ($\sigma_{\text{folds}}$):} Mide la estabilidad del modelo; valores elevados indican sobreajuste al conjunto de entrenamiento y guían la selección final de modelos.
    \item \textbf{Brecha público-local ($\Delta_{\text{lb}}$):} Diferencia entre el score en el \emph{leaderboard} público y el score en validación local; una brecha amplia señala sobreajuste al leaderboard ruidoso.
\end{itemize}

% -----------------------------------------------------------
\subsection{Validación cruzada temporal}

Dado que el conjunto de prueba corresponde a un periodo estrictamente posterior al entrenamiento, se empleará una estrategia de \emph{expanding window cross-validation}: cada fold extiende la ventana de entrenamiento hacia el futuro y reserva el bloque temporal inmediatamente siguiente como validación, simulando la predicción fuera de muestra. La selección final de modelos priorizará la estabilidad entre folds (menor $\sigma_{\text{folds}}$) sobre el score promedio máximo, para mitigar el sobreajuste al leaderboard público.


% -----------------------------------------------------------
\subsection{Ensemble por mezcla ponderada de rangos}

Como etapa final del pipeline, las predicciones de $K$ modelos base se combinan mediante un algoritmo de \emph{rank-based blending} con pesos diferenciados por posición. Para cada observación, los modelos se ordenan por magnitud de predicción en dos sentidos (descendente y ascendente), asignando a cada posición de rango $i$ un peso adicional $p_i$ sobre el peso base $w_k$ del modelo:

\begin{equation}
  \hat{y}_{\text{ens}} = w_{\downarrow} \cdot B_{\downarrow} + w_{\uparrow} \cdot B_{\uparrow}, \qquad
  B_{\downarrow} = \sum_{i=1}^{K} \hat{y}_{\sigma_i} \cdot \bigl(w_{\sigma_i} + p_i\bigr)
  \label{eq:ensemble}
\end{equation}

donde $\sigma = \arg\mathrm{sort}_{\downarrow}(\hat{\mathbf{y}})$ y $B_{\uparrow}$ se define análogamente con orden ascendente. Los pesos globales son $w_{\downarrow} = 0.70$ y $w_{\uparrow} = 0.30$; los pesos de posición son $\mathbf{p} = [+0.12, +0.06, +0.03, -0.01, -0.04, -0.06, -0.10]$. Este esquema penaliza las predicciones extremas de los modelos de menor peso y amplifica la señal del modelo dominante.


% ============================================================
% 6. EXPERIMENTOS Y RESULTADOS
% ============================================================

\section{Experimentos y Resultados}

% -----------------------------------------------------------
\subsection{Experimento 1: LightGBM multi-horizonte con ingeniería de features causales (Score: 0.2675)}

\paragraph{Datos y validación temporal.}
El conjunto de entrenamiento comprende 5\,337\,414 observaciones con $\texttt{ts\_index} \in [1, 3601]$ y 94 columnas. El conjunto de prueba contiene 1\,447\,107 observaciones con $\texttt{ts\_index} \in [3602, 4376]$ (92 columnas, sin target). La validación local usa la partición temporal $\texttt{ts\_index} > 3500$ como holdout, simulando predicción fuera de muestra sin filtración. El dataset cuenta con 23 códigos únicos, 180 sub-códigos y 5 sub-categorías; los pesos están fuertemente concentrados (el 10\% superior de las filas concentra el 98.3\% de la masa de peso total).

\paragraph{Estadísticas del objetivo por horizonte.}
La Tabla~\ref{tab:target_stats} muestra la alta varianza y curtosis del objetivo, especialmente en horizontes cortos, evidenciando el bajo SNR del problema.

\begin{table}[h!]
  \caption{Estadísticas descriptivas de \texttt{y\_target} por horizonte (conjunto de entrenamiento).}
  \label{tab:target_stats}
  \centering
  \begin{tabular}{ccccc}
    \toprule
    Horizonte & Media & Desv. estándar & Asimetría & Curtosis \\
    \midrule
    1  & $-0.083$ & $11.700$ & $ 5.28$ & $528.5$ \\
    3  & $-0.252$ & $19.361$ & $ 2.01$ & $243.8$ \\
    10 & $-0.776$ & $33.842$ & $ 0.83$ & $176.8$ \\
    25 & $-1.682$ & $52.823$ & $ 0.85$ & $141.5$ \\
    \bottomrule
  \end{tabular}
\end{table}

\paragraph{Ingeniería de features causales (183 features).}
Todos los estadísticos se calculan condicionados a $\texttt{ts\_index} \leq t$. Los grupos de features construidos se resumen en la Tabla~\ref{tab:feat_groups}.

\begin{table}[h!]
  \caption{Grupos de features construidos; todos son causales ($\texttt{ts\_index} \leq t$).}
  \label{tab:feat_groups}
  \centering
  \begin{tabular}{lp{7.5cm}}
    \toprule
    Grupo & Descripción \\
    \midrule
    Rezagos            & \texttt{lag\_1, lag\_3, lag\_7, lag\_14} sobre \texttt{feature\_al}; ratios entre rezagos (\texttt{lag\_ratio\_1\_7, \_3\_14, \_1\_3}). \\
    Rolling            & Media y desviación estándar con ventanas de 3, 7 y 14 pasos (shift(1) para evitar fuga); \texttt{trend\_strength}, \texttt{volatility\_strength}. \\
    EWM                & Medias exponenciales con spans $\{3,5\}$ para $h \leq 3$ y $\{5,7\}$ para $h > 3$ sobre \texttt{feature\_al, feature\_am, spread\_al\_am}. \\
    Estacionalidad     & \texttt{ts\_mod\_7/30/90}, \texttt{ts\_sin/cos} (período 365), \texttt{ts\_sin\_100/cos\_100} (período 100), \texttt{ts\_log}, \texttt{ts\_sq}. \\
    Interacción        & \texttt{horizon\_sq}, \texttt{ts\_horizon} (producto $\texttt{ts\_index} \times h$), \texttt{code\_horizon\_freq}. \\
    Jerárquico         & Media y desviación estándar de \texttt{y\_target} por \texttt{code}, \texttt{sub\_code} y \texttt{sub\_category} (calculados sobre $\texttt{ts\_index} \leq 3500$). \\
    Cross-seccional    & Z-score, rango percentil, min/max/rango y volatilidad normalizada de features clave (\texttt{feature\_al, am, cg, by, s, t}) por \texttt{ts\_index}. \\
    Spreads/Ratios     & \texttt{spread\_al\_am}, \texttt{ratio\_al\_am}, \texttt{spread\_cg\_by}, \texttt{ratio\_cg\_by}, \texttt{spread\_s\_t}, \texttt{ratio\_s\_t}; diferencias y variaciones porcentuales. \\
    Frecuencia         & \texttt{sub\_code\_freq}, \texttt{sub\_code\_log\_freq}, \texttt{sub\_code\_freq\_rank}, frecuencia relativa por timestamp. \\
    Dummies            & Variables indicadoras de \texttt{sub\_category} (5 categorías). \\
    \bottomrule
  \end{tabular}
\end{table}

\paragraph{Hiperparámetros de LightGBM por horizonte.}
Se entrena un modelo independiente para cada horizonte con la regularización ajustada al nivel de ruido: mayor regularización en $h=1$ (SNR más bajo) y menor en $h=25$. La Tabla~\ref{tab:lgb_params} resume la configuración.

\begin{table}[h!]
  \caption{Hiperparámetros de LightGBM por horizonte. Parámetros globales: \texttt{learning\_rate}=0.015, \texttt{n\_estimators}=4000, \texttt{feature\_fraction}=0.6, \texttt{bagging\_fraction}=0.7, \texttt{bagging\_freq}=5, $\lambda_1$=0.1, objetivo=RMSE.}
  \label{tab:lgb_params}
  \centering
  \begin{tabular}{ccccc}
    \toprule
    Horizonte & Hojas & Min. muestras/hoja & $\lambda_2$ & Early stopping \\
    \midrule
    1  & 70 & 250 & 12.0 & 200 \\
    3  & 75 & 225 & 11.0 & 190 \\
    10 & 85 & 180 &  9.0 & 170 \\
    25 & 90 & 150 &  8.0 & 160 \\
    \bottomrule
  \end{tabular}
\end{table}

El ensemble multi-semilla usa 5 semillas (42, 2024, 777, 1337, 9999); las predicciones de cada semilla se promedian antes de concatenar los horizontes.

\paragraph{Scores locales por horizonte y score global.}
La Tabla~\ref{tab:horizon_scores} presenta los scores de validación local (\emph{Weighted RMSE Skill Score}) para cada horizonte y el score agregado sobre el holdout completo.

\begin{table}[h!]
  \caption{Score de validación local por horizonte (holdout $\texttt{ts\_index} > 3500$). El score global se calcula sobre todos los horizontes simultáneamente.}
  \label{tab:horizon_scores}
  \centering
  \begin{tabular}{ccccc}
    \toprule
    Horizonte & Filas entrenamiento & Filas validación & Score local \\
    \midrule
    1  & 1\,351\,193 & 43\,460 & 0.07547 \\
    3  & 1\,342\,793 & 43\,023 & 0.13232 \\
    10 & 1\,296\,269 & 40\,967 & 0.23661 \\
    25 & 1\,181\,897 & 37\,812 & 0.29283 \\
    \midrule
    \textbf{Global} & --- & --- & \textbf{0.25083} \\
    \bottomrule
  \end{tabular}
\end{table}

El score crece con el horizonte porque las series de mayor duración tienen varianzas más altas (Tabla~\ref{tab:target_stats}), facilitando la extracción de señal por el modelo.

\paragraph{Importancia de features por horizonte.}
Para $h=1$, las features más influyentes son estadísticos cross-seccionales (\texttt{feature\_am\_ts\_min}, \texttt{feature\_al\_norm\_vol}, \texttt{sub\_cat\_std}) y features de estacionalidad de período corto (\texttt{ts\_cos\_100}, \texttt{ts\_mod\_90}). Para $h \geq 10$, las features de tiempo dominan (\texttt{ts\_log}, \texttt{ts\_cos}, \texttt{ts\_sin}) junto con estadísticos jerárquicos (\texttt{sub\_cat\_mean}, \texttt{sub\_cat\_std}), evidenciando que los horizontes largos capturan tendencias de largo plazo.

\paragraph{Estadísticas del vector de predicción final y score público.}
El modelo generó 1\,447\,107 predicciones con media $-0.636$, desviación estándar $2.728$, mínimo $-62.98$ y máximo $12.81$. El \textbf{score en el leaderboard público fue 0.2675}, superando al mejor modelo individual del experimento de ensemble (0.2638) en $\Delta = +0.0037$.

% -----------------------------------------------------------
\subsection{Experimento 3: LightGBM con Polars, suavizado de outliers y selección de features (Score: 0.2712)}

Este experimento introduce tres innovaciones sobre el Experimento 1: (i)~un pipeline de preprocesamiento en \textbf{Polars} con reducción de memoria, (ii)~\textbf{suavizado de outliers} tipo Hampel, y (iii)~\textbf{ingeniería de interacciones} con selección automática de features por horizonte. El entrenamiento se paraleliza entre horizontes mediante \texttt{joblib}.

\paragraph{Preprocesamiento con Polars.}
Los datos se cargan directamente en Polars por horizonte y se someten a tres etapas antes de llegar a LightGBM:

\begin{enumerate}
    \item \textbf{Forward-fill causal combinado}: train y test se concatenan (marcados con \texttt{\_is\_test}), se ordenan por grupo y \texttt{ts\_index}, y se aplica \texttt{forward\_fill} sobre todas las features numéricas dentro de cada grupo. Esto permite al test \emph{tomar prestada} la última observación histórica del train sin filtrar el target, que se excluye explícitamente. Nulls residuales al inicio de cada serie se rellenan con cero.
    \item \textbf{Reducción de memoria}: Float64 → Float32, Int64 → Int8/16/32 (según rango), strings con cardinalidad baja → Categorical. La memoria se reduce un \textbf{47\%} (de $\approx$1\,023\,MB a $\approx$540\,MB por horizonte).
    \item \textbf{Suavizado Pseudo-Hampel}: sobre 12 features ruidosas (\texttt{feature\_al}, \texttt{feature\_x}, \texttt{feature\_n}, \texttt{feature\_y}, \texttt{feature\_z}, \texttt{feature\_bo}, \texttt{feature\_bm}, \texttt{feature\_bz}, \texttt{feature\_ag}, \texttt{feature\_ao}, \texttt{feature\_af}, \texttt{feature\_u}). Para cada feature, se calcula la mediana rolling ($w=50$) y la desviación estándar rolling ($w=50$) dentro de cada grupo; los valores que superan $5\sigma$ respecto a la mediana local son reemplazados por esa mediana.
\end{enumerate}

\paragraph{Ingeniería de features avanzada.}
Sobre los datos preprocesados se construyen las siguientes familias de features, todas causales:

\begin{itemize}
    \item \textbf{Features de signo y magnitud}: variables categóricas del signo de \texttt{feature\_cd} y \texttt{feature\_bz} ($\in \{-1, 0, +1\}$); y productos $\mathrm{sign}(A) \times |B|$ para dos pares de features (\texttt{feature\_cd}$\times$\texttt{feature\_af}; \texttt{feature\_v}$\times$\texttt{feature\_i}).
    \item \textbf{Features de interacción}: para 9 pares de features seleccionadas por correlación con el target, se generan 5 operaciones por par (suma, diferencia, producto, ratio $A/B$ y ratio $B/A$), produciendo hasta 45 features de interacción.
    \item \textbf{Features rolling acumulativas} (ventana $w=1000$): mediana rolling, máximo rolling, mínimo rolling y Z-score robusto ($(x - \text{mediana}_{w}) / (\sigma_w + \epsilon)$) sobre 4 features (\texttt{feature\_bz}, \texttt{feature\_bo}, \texttt{feature\_v}, \texttt{feature\_i}).
    \item \textbf{Features de momentum acumulativo}: expanding median, EMA con spans $\{10, 30\}$, diferencia respecto al EMA, ratio respecto al EMA, MACD (EMA$_{10}$ $-$ EMA$_{30}$) y \emph{trend conflict} (EMA$_{10}$ $-$ expanding median); calculadas sobre \texttt{inter\_feature\_am\_mul\_feature\_cc} y \texttt{feature\_al}.
    \item \textbf{Features row-wise}: máximo, mínimo y amplitud horizontal para dos grupos de tres features (\texttt{feature\_af}/\texttt{feature\_u}/\texttt{feature\_bm} y \texttt{feature\_bz}/\texttt{feature\_cc}/\texttt{feature\_by}).
    \item \textbf{Features clásicas} (Pandas, al convertir a DataFrame): rezagos 1/3/5/10/25 sobre 4 features, rolling mean/std con ventanas 5/10/20 sobre 2 features, pendiente rolling (\texttt{polyfit}) con ventanas 10/20, diferencia de primer orden, rango percentil cross-seccional, ciclo temporal $\sin(2\pi \cdot t / 100)$, y codificación de \texttt{sub\_category}/\texttt{sub\_code} por media de target.
    \item \textbf{Eliminación de 17 features} de baja correlación con el target (\texttt{feature\_bb}, \texttt{feature\_t}, \texttt{feature\_o}, \texttt{feature\_s}, \texttt{feature\_ba}, \texttt{feature\_bc}, \texttt{feature\_p}, \texttt{feature\_aw}, \texttt{feature\_bm}, \texttt{feature\_cc}, \texttt{feature\_an}, \texttt{feature\_bo}, \texttt{feature\_u}, \texttt{feature\_bz}, \texttt{feature\_ap}, \texttt{feature\_ae}, \texttt{feature\_q}).
\end{itemize}

\paragraph{Selección de features en dos pasadas.}
Debido al alto número de features generadas, se aplica una selección en dos pasadas por horizonte:
\begin{enumerate}
    \item \textbf{Pasada de ranking}: se entrena un LightGBM rápido (early stopping en 50 rondas, sin pesos) y se ordenan todas las features por \emph{gain} de importancia.
    \item \textbf{Pasada final}: se retienen únicamente las $N_h$ features más importantes según el objetivo por horizonte: $N_1 = 197$, $N_3 = 177$, $N_{10} = 197$, $N_{25} = 217$.
\end{enumerate}

\paragraph{Configuración de LightGBM.}
A diferencia del Experimento 1, este experimento usa un único conjunto de hiperparámetros para todos los horizontes: \texttt{num\_leaves}=31, \texttt{min\_child\_samples}=200, $\lambda_2=10$, $\lambda_1=0.1$, \texttt{learning\_rate}=0.015, \texttt{n\_estimators}=4200, \texttt{feature\_fraction}=0.6, \texttt{bagging\_fraction}=0.7, \texttt{bagging\_freq}=5. El ensemble usa 5 semillas [42, 2026, 1\,232\,006, 3\,107, 2\,709]. Los cuatro horizontes se entrenan en paralelo (\texttt{joblib}, \texttt{n\_jobs}=2).

\paragraph{Scores locales por horizonte y score público.}
La Tabla~\ref{tab:horizon_scores_exp3} muestra los scores de validación local (holdout $\texttt{ts\_index} > 3500$) para cada horizonte.

\begin{table}[h!]
  \caption{Score de validación local por horizonte — Experimento 3. El score global en el leaderboard público fue 0.2712.}
  \label{tab:horizon_scores_exp3}
  \centering
  \begin{tabular}{ccc}
    \toprule
    Horizonte & Features seleccionadas ($N_h$) & Score local \\
    \midrule
    1  & 197 & 0.08037 \\
    3  & 177 & 0.11295 \\
    10 & 197 & 0.20699 \\
    25 & 217 & 0.28280 \\
    \bottomrule
  \end{tabular}
\end{table}

El \textbf{score en el leaderboard público fue 0.2712}, superando al Experimento 1 (0.2675) en $\Delta = +0.0037$. La discrepancia entre los scores locales por horizonte (menores que los del Experimento 1) y el mejor score público evidencia que la validación local single-window no captura perfectamente la distribución del conjunto de prueba privado, lo que refuerza la importancia del criterio de estabilidad entre folds establecido en los objetivos.

% -----------------------------------------------------------
\subsection{Experimento 4: Stacking LightGBM + XGBoost + Ridge con pesos por recencia (Score: 0.2824)}

Este experimento introduce dos innovaciones estructurales sobre los anteriores: (i)~un \textbf{meta-modelo de stacking} que combina LightGBM y XGBoost con una Ridge Regression entrenada fuera de muestra respetando el orden temporal, y (ii)~una \textbf{ponderación por recencia} que amplifica el peso de las observaciones más recientes durante el entrenamiento. Es el experimento con el mayor score público obtenido (0.2824).

\paragraph{Preprocesamiento y construcción del conjunto de entrenamiento.}
Por cada horizonte se lee solo el subconjunto correspondiente del parquet. Para calcular features de rezago sobre el test sin filtrar el futuro, se anteponen las últimas 50 observaciones del train a cada lote de test (\emph{history injection}); las predicciones se extraen únicamente sobre las filas con $\texttt{ts\_index} > \max(\texttt{ts\_index}_{\text{train}})$.

El umbral de validación se ajusta por horizonte para evitar que los rezagos más largos contaminen el holdout: $t_{\text{fit}} = 3500 - (h + 25)$ para $h \in \{1, 10\}$, $t_{\text{fit}} = 3500 - 25$ para $h = 3$, y $t_{\text{fit}} = 3500$ para $h = 25$.

\paragraph{Ponderación por recencia.}
Los pesos de entrenamiento originales ($w_i$) se amplifican mediante una función lineal del índice temporal normalizado:
\begin{equation}
  \tilde{w}_i = w_i \cdot \bigl(0.5 + 40.0 \cdot \hat{t}_i\bigr), \qquad
  \hat{t}_i = \frac{t_i - t_{\min}}{t_{\max} - t_{\min} + \varepsilon}
  \label{eq:recency}
\end{equation}
donde $\hat{t}_i \in [0,1]$ es el índice temporal normalizado. Las observaciones más recientes ($\hat{t} \approx 1$) reciben hasta $40.5\times$ más peso que las más antiguas ($\hat{t} \approx 0$), forzando al modelo a ajustarse a la dinámica reciente del mercado.

\paragraph{Ingeniería de features híbrida por horizonte.}
Se construyen features sobre 5 columnas clave (\texttt{feature\_al}, \texttt{feature\_am}, \texttt{feature\_cg}, \texttt{feature\_by}, \texttt{feature\_s}). La Tabla~\ref{tab:hybrid_feats} resume los grupos de features comunes y los específicos por horizonte.

\begin{table}[h!]
  \caption{Grupos de features de la ingeniería híbrida. Todas son causales; las específicas por horizonte se aplican adicionalmente sobre los grupos comunes.}
  \label{tab:hybrid_feats}
  \centering
  \begin{tabular}{lp{9cm}}
    \toprule
    Grupo & Descripción \\
    \midrule
    Frecuencia & \texttt{sub\_code\_freq}, log-freq, rank, frecuencia relativa por \texttt{ts\_index}. \\
    Diferenciales & \texttt{d\_al\_am}, \texttt{r\_al\_am}, \texttt{d\_cg\_by}, \texttt{d\_u\_s}, \texttt{d\_ac\_ab}, \texttt{d\_m\_l}. \\
    Norm. cross-secc. & Z-score crudo de 5 features por \texttt{ts\_index} (\texttt{\_cs}); contador de outliers ($|\text{cs}| > 3$). \\
    Rezagos & \texttt{lag}\{1,3,5,10,25\} para cada una de las 5 features. \\
    EMA + Wavelet & EMA con spans \{3,10,25\}; diferencias EMA$_3$-EMA$_{10}$ y EMA$_{10}$-EMA$_{25}$ como proxy de momentum multi-escala (\emph{wavelet}). \\
    Régimen de tendencia & Señal ternaria $\{-1, 0, +1\}$ según el alineamiento de ambas wavelets. \\
    Oscilador estocástico & $(\texttt{x} - \min_{25}) / (\max_{25} - \min_{25})$ (ventana 25). \\
    Reversión a la media & $(\texttt{x} - \text{EMA}_3) / (\sigma_{3} + \varepsilon)$; ratio de volatilidades $\sigma_3 / \sigma_{10}$ (\emph{vol squeeze}). \\
    Vs. sub-código & Desviación de cada instrumento respecto a la media de su \texttt{sub\_code} en el mismo \texttt{ts\_index}. \\
    Tendencia lineal & Pendiente \texttt{polyfit}(10) para \texttt{feature\_al} y \texttt{feature\_am}. \\
    \midrule
    $h = 25$ solamente & \texttt{lag50}, EMA$_{50}$, wavelet larga (EMA$_{25}$-EMA$_{50}$), \texttt{trend\_25}. \\
    $h \in \{3,10\}$ solamente & Cross-seccional de diff1 y wavelets; outlier count global (solo $h=10$). \\
    $h \in \{1,3\}$ solamente & Decil de momentum de \texttt{feature\_cg}, divergencia \texttt{cg}$\times$\texttt{al\_cs}, rank de reversión. \\
    $h = 1$ solamente & Conteo de missings por fila. \\
    \bottomrule
  \end{tabular}
\end{table}

\paragraph{Configuración de LightGBM y XGBoost.}
La Tabla~\ref{tab:lgb_xgb_cfg} muestra los hiperparámetros principales de ambos modelos, entrenados con las mismas 5 semillas [42, 2024, 12345, 666, 7777] y promediados antes del stacking. Las predicciones de cada semilla se acumulan con peso $1/5$.

\begin{table}[h!]
  \caption{Hiperparámetros de LightGBM y XGBoost. Ambos usan early stopping; LGB en 100 rondas, XGB en 300. Los parámetros son uniformes entre horizontes.}
  \label{tab:lgb_xgb_cfg}
  \centering
  \begin{tabular}{lcc}
    \toprule
    Parámetro & LightGBM & XGBoost \\
    \midrule
    \texttt{learning\_rate}      & 0.005 & 0.015 \\
    \texttt{n\_estimators}       & 6\,000 & 6\,000 \\
    \texttt{num\_leaves / max\_leaves} & 1\,024 & 1\,024 \\
    \texttt{max\_depth}          & 14 & 10 \\
    \texttt{min\_child\_samples / min\_child\_weight} & 400 & 5.0 \\
    \texttt{feature\_fraction / colsample\_bytree} & 0.6 & 0.6 \\
    \texttt{bagging\_fraction / subsample} & 0.8 & 0.8 \\
    $\lambda_1$ (\texttt{reg\_alpha})   & 0.1 & 0.1 \\
    $\lambda_2$ (\texttt{reg\_lambda})  & 10.0 & 1.0 \\
    \texttt{grow\_policy}        & leaf-wise & \texttt{lossguide} \\
    Datos categóricos            & --- & \texttt{enable\_categorical=True} \\
    \bottomrule
  \end{tabular}
\end{table}

\paragraph{Meta-modelo Ridge con validación temporal.}
Las predicciones de validación del ensemble LGB ($\hat{y}_{\text{LGB}}$) y XGB ($\hat{y}_{\text{XGB}}$) se apilan como columnas de una matriz $\mathbf{M} \in \mathbb{R}^{n_{\text{val}} \times 2}$. Sobre esta matriz se entrena una regresión Ridge ($\alpha = 1$, pesos no negativos, \texttt{positive=True}) usando 5-Fold KFold \emph{sin barajado} (preservando el orden temporal) para generar predicciones fuera de muestra:
\begin{equation}
  \hat{y}_{\text{stack}} = w_{\text{LGB}} \cdot \hat{y}_{\text{LGB}} + w_{\text{XGB}} \cdot \hat{y}_{\text{XGB}}, \qquad w_{\text{LGB}}, w_{\text{XGB}} \geq 0
  \label{eq:ridge_stack}
\end{equation}
Para la predicción final en test se ajusta un Ridge sobre el holdout completo.

\paragraph{Resultados locales por horizonte y score público.}
La Tabla~\ref{tab:stacking_scores} detalla los scores de validación local para cada modelo y horizonte, junto con los pesos aprendidos por el meta-modelo.

\begin{table}[h!]
  \caption{Scores locales de validación ($\texttt{ts\_index} > 3500$) por horizonte y modelo. Los pesos Ridge suman 100\% por horizonte. Score global de stacking: 0.24446. Score en el leaderboard público: \textbf{0.2824}.}
  \label{tab:stacking_scores}
  \centering
  \begin{tabular}{ccccccc}
    \toprule
    $h$ & LGB local & XGB local & Stacking local & $w_{\text{LGB}}$ & $w_{\text{XGB}}$ \\
    \midrule
    1  & 0.08947 & 0.01490 & 0.08109 & 78.6\% & 21.4\% \\
    3  & 0.13313 & 0.03721 & 0.13573 & 64.4\% & 35.6\% \\
    10 & 0.24442 & 0.03942 & 0.24379 & 94.3\% &  5.7\% \\
    25 & 0.29130 & 0.15028 & 0.27791 & 90.6\% &  9.4\% \\
    \midrule
    Global & --- & --- & \textbf{0.24446} & --- & --- \\
    \bottomrule
  \end{tabular}
\end{table}

Tres patrones destacan de la Tabla~\ref{tab:stacking_scores}: (1)~LightGBM domina sistemáticamente sobre XGBoost en este problema (pesos Ridge 64--94\%), sugiriendo que la estrategia de crecimiento hoja-a-hoja de LightGBM captura mejor la señal de baja SNR; (2)~el stacking mejora el score individual de LGB únicamente en $h=3$ ($+0.003$), mientras que en los demás horizontes lo iguala o reduce marginalmente, indicando que XGBoost aporta señal complementaria principalmente en horizontes intermedios; (3)~el score de validación local del stacking (0.24446) es inferior al del Experimento 1 (0.25083), pero el \textbf{score en el leaderboard público (0.2824) supera en $\Delta = +0.0149$ al mejor resultado previo} (Experimento 3, 0.2675), evidenciando que la ponderación por recencia y las features de régimen capturan dinámicas del conjunto de prueba privado que el holdout local no refleja.

% -----------------------------------------------------------
\subsection{Experimento 5: LightGBM minimalista con identificadores jerárquicos como features (Score: 0.3212)}

Este experimento es el resultado más importante del trabajo y el que mejor encarna el título \emph{No Free Lunch}: el modelo más simple de toda la serie —sin ensemble de semillas, sin ingeniería de features elaborada, sin apilamiento— obtiene el mayor score público (0.3212), superando al stacking de cuatro modelos del Experimento 4 en $\Delta = +0.0388$.

\paragraph{Diseño del modelo.}
El pipeline es intencionalmente minimalista:
\begin{enumerate}
    \item Se leen los 86 features originales más las columnas de identidad (\texttt{id}, \texttt{code}, \texttt{sub\_code}, \texttt{sub\_category}, \texttt{horizon}).
    \item Se añaden únicamente tres features derivadas: la diferencia $\texttt{feature\_al} - \texttt{feature\_am}$, y la media global por grupo $(\texttt{code}, \texttt{sub\_code}, \texttt{sub\_category}, \texttt{horizon})$ de \texttt{feature\_al} y \texttt{feature\_am} respectivamente.\footnote{Las medias de grupo se calculan sobre todo el split correspondiente (train o test), lo que introduce una fuente menor de información futura en las features —aunque no en el target. Se reporta como limitación del experimento.}
    \item Todas las columnas de tipo \texttt{object} (\texttt{id}, \texttt{code}, \texttt{sub\_code}, \texttt{sub\_category}) se convierten a \texttt{category} y se pasan directamente a LightGBM con \texttt{categorical\_feature='auto'}.
    \item Se entrena un único modelo por horizonte, sin ensemble de semillas ni meta-modelo.
\end{enumerate}

El uso de los identificadores jerárquicos como features categóricas le permite a LightGBM aprender \textbf{efectos fijos por entidad}: cada combinación única de (\texttt{code}, \texttt{sub\_code}, \texttt{sub\_category}) recibe implícitamente un intercepto específico, capturando los niveles base de cada instrumento sin necesidad de codificación explícita. La columna \texttt{id} actúa como clave única por observación de entrenamiento (cada \texttt{ts\_index} produce un \texttt{id} distinto); para el test, los \texttt{id} son todos nuevos (horizonte temporal desplazado), por lo que LightGBM enruta dichas observaciones por las ramas de las demás features, eliminando el riesgo de memorización pura.

\paragraph{Hiperparámetros.}
La configuración es notablemente más simple que en los experimentos anteriores: \texttt{learning\_rate}=0.035 (el más alto de la serie), \texttt{n\_estimators}=1202 con early stopping en 100 rondas, \texttt{num\_leaves}=64, \texttt{min\_child\_samples}=100, \texttt{feature\_fraction}=0.75. No se especifican parámetros de bagging ni regularización $\lambda$.

\paragraph{Validación y métrica local.}
La validación usa el corte temporal estándar ($\texttt{ts\_index} > 3500$). La función de score local en este experimento difiere de la métrica de competencia: calcula $\sqrt{\text{WSSE}/\text{WSSY}^2}$ (raíz de la razón, \emph{sin} el complemento $1-$), por lo que sus valores son menores-es-mejor y no son directamente comparables con los scores de los experimentos anteriores. La Tabla~\ref{tab:exp5_conv} muestra la convergencia del RMSE de validación (no ponderado) para cada horizonte.

\begin{table}[h!]
  \caption{Convergencia por horizonte — Experimento 5. RMSE de validación es la métrica interna de LightGBM (no ponderado, sin el complemento $1-$). El score público fue \textbf{0.3212}.}
  \label{tab:exp5_conv}
  \centering
  \begin{tabular}{ccc}
    \toprule
    Horizonte & Mejor iteración & RMSE validación \\
    \midrule
    1  & 133 & 0.001134 \\
    3  &  93 & 0.001818 \\
    10 & 107 & 0.002676 \\
    25 & 117 & 0.002908 \\
    \bottomrule
  \end{tabular}
\end{table}

Los modelos convergen muy rápido (93--133 iteraciones de 1202 posibles), lo que indica que la señal útil se extrae en pocas rondas y el early stopping actúa como regularizador efectivo.

\paragraph{Discusión: la paradoja de la simplicidad.}
La Tabla~\ref{tab:resumen_scores} compara los scores públicos de todos los experimentos. El Experimento 5 supera a todos pese a ser el más simple, lo que constituye una demostración directa del teorema \emph{No Free Lunch}: no existe un método uniformemente superior en todos los dominios. En este problema de baja SNR con series anónimas, agregar features elaboradas, ensembles multi-semilla y meta-modelos introduce ruido que degrada la generalización fuera de muestra. Los efectos fijos por entidad —capturados de forma gratuita por LightGBM al tratar \texttt{code}, \texttt{sub\_code} y \texttt{sub\_category} como categorías— contienen más señal predictiva que cientos de estadísticos rolling.

\begin{table}[h!]
  \caption{Comparación de scores públicos en el leaderboard entre todos los experimentos.}
  \label{tab:resumen_scores}
  \centering
  \begin{tabular}{clc}
    \toprule
    Experimento & Descripción & Score público \\
    \midrule
    2 & Ensemble rank-blending (7 modelos)                          & 0.2639 \\
    1 & LightGBM multi-horizonte, 183 features, 5 semillas          & 0.2675 \\
    3 & LightGBM + Polars + Hampel + interacciones                  & 0.2712 \\
    4 & LGB + XGB + Ridge stacking, recency weighting               & 0.2824 \\
    \textbf{5} & \textbf{LightGBM minimalista, IDs jerárquicos como features} & \textbf{0.3212} \\
    \bottomrule
  \end{tabular}
\end{table}

% -----------------------------------------------------------
\subsection{Experimento 6: LightGBM minimalista con umbral de validación ajustado y entrenamiento determinista (Score: 0.3377)}

El Experimento 6 es una refinación directa del Experimento 5 con el mismo pipeline minimalista, pero con tres ajustes quirúrgicos que elevan el score de 0.3212 a \textbf{0.3377} ($\Delta = +0.0165$), constituyendo el mejor resultado de toda la serie experimental.

\paragraph{Cambios respecto al Experimento 5.}
La Tabla~\ref{tab:diff_exp5_exp6} sintetiza las diferencias entre ambos experimentos. El resto del pipeline —features, arquitectura de modelo, early stopping en 100 rondas, métrica local— es idéntico.

\begin{table}[h!]
  \caption{Diferencias entre el Experimento 5 y el Experimento 6. Todo lo no listado es idéntico.}
  \label{tab:diff_exp5_exp6}
  \centering
  \begin{tabular}{lcc}
    \toprule
    Parámetro & Experimento 5 & Experimento 6 \\
    \midrule
    \texttt{VAL\_THRESHOLD}     & 3\,500 & \textbf{3\,503} \\
    \texttt{learning\_rate}     & 0.035  & \textbf{0.030}  \\
    \texttt{feature\_fraction}  & 0.75   & \textbf{0.70}   \\
    Semilla                     & por defecto & \textbf{\texttt{seed=5}} \\
    \texttt{deterministic}      & \texttt{False} & \textbf{\texttt{True}} \\
    \texttt{num\_threads}       & por defecto & \textbf{1} \\
    Score público               & 0.3212 & \textbf{0.3377} \\
    \bottomrule
  \end{tabular}
\end{table}

\paragraph{Efecto del umbral de validación.}
El desplazamiento de \texttt{VAL\_THRESHOLD} de 3500 a 3503 incorpora tres pasos temporales adicionales al conjunto de entrenamiento. El impacto desproporcionado en el score ($+0.0165$) sugiere que esos tres índices contienen observaciones de alta información para la dinámica del modelo, o que el punto de corte 3503 produce una partición validación/test con distribución más alineada a la del test privado. Este hallazgo evidencia la alta sensibilidad del pipeline al punto de validación en series de baja SNR.

\paragraph{Entrenamiento determinista.}
La combinación de \texttt{seed=5}, \texttt{deterministic=True}, \texttt{force\_row\_wise=True} y \texttt{num\_threads=1} garantiza resultados bit-a-bit reproducibles entre ejecuciones, lo cual es crítico para identificar si mejoras marginales provienen de ajustes reales o de varianza estocástica del entrenamiento.

\paragraph{Convergencia por horizonte.}
La Tabla~\ref{tab:exp6_conv} muestra la convergencia del RMSE de validación (no ponderado). En comparación con el Experimento 5, los horizontes cortos convergen más rápido (h=1: 53 vs 133 iteraciones) mientras que los horizontes largos requieren significativamente más iteraciones (h=10: 541 vs 107; h=25: 433 vs 117), lo que indica que la menor tasa de aprendizaje y la fracción de features más restrictiva favorecen una exploración más profunda del espacio de modelos para las series de largo plazo.

\begin{table}[h!]
  \caption{Convergencia por horizonte — Experimento 6. Score público: \textbf{0.3377}.}
  \label{tab:exp6_conv}
  \centering
  \begin{tabular}{cccc}
    \toprule
    Horizonte & Mejor iteración & RMSE validación & Mejor iter. (Exp.~5) \\
    \midrule
    1  &  53 & 0.001142 & 133 \\
    3  & 118 & 0.001825 &  93 \\
    10 & 541 & 0.002653 & 107 \\
    25 & 433 & 0.002843 & 117 \\
    \bottomrule
  \end{tabular}
\end{table}

\paragraph{Tabla comparativa final de todos los experimentos.}
La Tabla~\ref{tab:resumen_scores_final} actualiza la comparación incluyendo el Experimento 6.

\begin{table}[h!]
  \caption{Ranking final de scores públicos en el leaderboard. Los Experimentos 5 y 6 dominan con el diseño más simple.}
  \label{tab:resumen_scores_final}
  \centering
  \begin{tabular}{clcc}
    \toprule
    Exp. & Descripción & Score público & $\Delta$ vs. baseline \\
    \midrule
    2 & Ensemble rank-blending (7 modelos)                   & 0.2639 & --- \\
    1 & LightGBM, 183 features, 5 semillas                   & 0.2675 & $+0.0036$ \\
    3 & LightGBM + Polars + Hampel + interacciones           & 0.2712 & $+0.0073$ \\
    4 & LGB + XGB + Ridge stacking, recency weighting        & 0.2824 & $+0.0185$ \\
    5 & LightGBM minimalista, IDs jerárquicos, thr=3500      & 0.3212 & $+0.0573$ \\
    \textbf{6} & \textbf{LightGBM minimalista, IDs jerárquicos, thr=3503} & \textbf{0.3377} & $\mathbf{+0.0738}$ \\
    \bottomrule
  \end{tabular}
\end{table}

La progresión de scores revela que los saltos más grandes provienen de decisiones arquitectónicas (Exp.~4$\to$5: $+0.0388$) y no de hiperparámetros finos. Sin embargo, la afinación del umbral de validación (Exp.~5$\to$6: $+0.0165$) tiene un impacto comparable al del salto entre los experimentos más elaborados (Exp.~3$\to$4: $+0.0112$), subrayando la importancia crítica de la estrategia de validación en series de tiempo financieras de baja SNR.

% -----------------------------------------------------------
\subsection{Experimento 7: Ensemble dual LightGBM — modelos por horizonte y por sub-categoría con mezcla ponderada por potencia (Score: 0.3429)}

Este experimento introduce una segunda dimensión de especialización: además de los 4 modelos por horizonte (vista temporal), se entrenan 5 modelos por sub-categoría (vista estructural). Las 9 predicciones se fusionan mediante mezcla ponderada por potencia orientada por los scores de validación locales, logrando el \textbf{score público más alto de la serie experimental: 0.3429}.

\paragraph{Arquitectura de doble pool.}
Se entrenan dos conjuntos de modelos en paralelo:
\begin{itemize}
    \item \textbf{Pool A — modelos por horizonte} (4 modelos): cada modelo recibe todas las sub-categorías y aprende los patrones temporales específicos de cada horizonte $h \in \{1, 3, 10, 25\}$.
    \item \textbf{Pool B — modelos por sub-categoría} (5 modelos): cada modelo recibe todos los horizontes y aprende los patrones estructurales de cada sub-categoría $\in \{$\texttt{DPPUO5X2, NQ58FVQM, PHHHVYZI, PZ9S1Z4V, V8BKY1IV}$\}$.
\end{itemize}
Cada test se predice dos veces —una desde el pool temporal y otra desde el pool estructural— y las predicciones se fusionan con pesos derivados de los scores de validación locales.

\paragraph{Ingeniería de features (96 features).}
La función \texttt{create\_features} añade 10 features derivadas a las 86 originales, con \texttt{id}, \texttt{weight} y \texttt{ts\_index} explícitamente excluidos (a diferencia de los Experimentos 5 y 6 donde \texttt{id} permanecía como feature):
\begin{itemize}
    \item \textbf{Diferencia}: \texttt{feature\_al\_minus\_feature\_am}.
    \item \textbf{Medias de grupo} por (\texttt{code}, \texttt{sub\_code}, \texttt{sub\_category}, \texttt{horizon}): \texttt{feature\_al\_grp\_mean}, \texttt{feature\_am\_grp\_mean}.
    \item \textbf{Codificación de target} (\emph{target encoding}): \texttt{sub\_code\_target\_mean} — media de \texttt{y\_target} por \texttt{sub\_code} calculada exclusivamente sobre $\texttt{ts\_index} \leq 3500$; observaciones con \texttt{sub\_code} desconocido reciben la media global ($-0.6912$).
    \item \textbf{Rezagos}: \texttt{feature\_al\_lag1} y \texttt{feature\_am\_lag1} (rezago 1 por grupo \texttt{code}$\times$\texttt{horizon}), con NaN rellenados con cero.
    \item \textbf{Categoriales}: \texttt{code}, \texttt{sub\_code}, \texttt{sub\_category}, \texttt{horizon} convertidos a \texttt{category}.
\end{itemize}

\paragraph{Configuración de LightGBM.}
Parámetros uniformes para los 9 modelos: \texttt{learning\_rate}=0.03, \texttt{n\_estimators}=1200, \texttt{num\_leaves}=64, \texttt{min\_child\_samples}=100, \texttt{feature\_fraction}=0.75, \texttt{bagging\_fraction}=0.75, \texttt{bagging\_freq}=5, \texttt{random\_state}=101. Early stopping en 100 rondas sobre RMSE de validación ponderado.

\paragraph{Scores de validación por modelo.}
La Tabla~\ref{tab:exp7_scores} muestra los scores de validación local (métrica correcta de la competencia) para cada modelo de ambos pools.

\begin{table}[h!]
  \caption{Scores de validación local por modelo — Experimento 7. Los modelos de sub-categoría superan sistemáticamente a los de horizonte corto, motivando la mezcla asimétrica.}
  \label{tab:exp7_scores}
  \centering
  \begin{tabular}{llc}
    \toprule
    Pool & Modelo & Score local \\
    \midrule
    \multirow{5}{*}{Por horizonte}
      & $h=1$   & 0.0731 \\
      & $h=3$   & 0.1289 \\
      & $h=10$  & 0.2510 \\
      & $h=25$  & 0.3271 \\
      & \textit{Promedio} & \textit{0.1950} \\
    \midrule
    \multirow{6}{*}{Por sub-categoría}
      & \texttt{DPPUO5X2} & 0.2126 \\
      & \texttt{NQ58FVQM} & 0.2367 \\
      & \texttt{PHHHVYZI} & 0.2879 \\
      & \texttt{PZ9S1Z4V} & 0.3133 \\
      & \texttt{V8BKY1IV} & 0.2203 \\
      & \textit{Promedio} & \textit{0.2541} \\
    \bottomrule
  \end{tabular}
\end{table}

\paragraph{Mezcla ponderada por potencia.}
Para cada observación de test con horizonte $h$ y sub-categoría $s$, la predicción final se calcula como:
\begin{equation}
  \hat{y} = \frac{\text{score}_h^{p}}{\text{score}_h^{p} + \text{score}_s^{p}} \cdot \hat{y}_h
          + \frac{\text{score}_s^{p}}{\text{score}_h^{p} + \text{score}_s^{p}} \cdot \hat{y}_s, \qquad p = 2.5
  \label{eq:power_blend}
\end{equation}
El exponente $p=2.5$ amplifica agresivamente la diferencia entre modelos: por ejemplo, la observación con $h=1$ (score=0.073) mezclada con \texttt{PZ9S1Z4V} (score=0.313) resulta en pesos 2.6\% y 97.4\% respectivamente. En promedio, el pool de sub-categoría recibe el \textbf{67.1\%} del peso total y el pool de horizonte el \textbf{32.9\%}, reflejando que la vista estructural captura más señal que la vista temporal en este dataset.

\paragraph{Estadísticas de la predicción final y score público.}
Las 1\,447\,107 predicciones tienen media $-0.869$, desviación estándar $7.440$, mínimo $-158.05$ y máximo $86.75$. La mayor dispersión respecto a los experimentos anteriores (std 2.7–4.1) refleja la contribución de los modelos de sub-categoría, que modelan regímenes con varianzas más altas (ej. \texttt{PHHHVYZI} y \texttt{PZ9S1Z4V}).

El \textbf{score en el leaderboard público fue 0.3429}, superando al Experimento 6 (0.3377) en $\Delta = +0.0052$ y al Experimento 5 (0.3212) en $\Delta = +0.0217$.

\paragraph{Tabla comparativa final actualizada.}
La Tabla~\ref{tab:resumen_final_7} extiende la comparación a los 7 experimentos.

\begin{table}[h!]
  \caption{Ranking final de scores públicos incluyendo todos los experimentos. Los tres mejores corresponden a arquitecturas minimalistas con identificadores jerárquicos.}
  \label{tab:resumen_final_7}
  \centering
  \begin{tabular}{clcc}
    \toprule
    Exp. & Descripción & Score público & $\Delta$ vs. Exp.~2 \\
    \midrule
    2 & Ensemble rank-blending (7 modelos)                          & 0.2639 & --- \\
    1 & LightGBM, 183 features, 5 semillas                          & 0.2675 & $+0.0036$ \\
    3 & LightGBM + Polars + Hampel + interacciones                  & 0.2712 & $+0.0073$ \\
    4 & LGB + XGB + Ridge stacking, recency weighting               & 0.2824 & $+0.0185$ \\
    5 & LightGBM minimalista, IDs jerárquicos, thr=3500             & 0.3212 & $+0.0573$ \\
    6 & LightGBM minimalista, IDs jerárquicos, thr=3503             & 0.3377 & $+0.0738$ \\
    \textbf{7} & \textbf{Dual ensemble: 4 horizonte + 5 sub-cat, mezcla $p^{2.5}$} & \textbf{0.3429} & $\mathbf{+0.0790}$ \\
    \bottomrule
  \end{tabular}
\end{table}

Los resultados confirman que la segmentación por sub-categoría es complementaria a la segmentación por horizonte: los modelos estructurales capturan diferencias de régimen entre grupos de instrumentos que los modelos temporales no pueden resolver, y la mezcla ponderada por potencia asigna los pesos de forma automática y orientada a datos.

% -----------------------------------------------------------
\subsection{Experimento 2: Ensemble por mezcla ponderada de rangos (Score: 0.2639)}

Como experimento complementario, se combinaron 7 submissions previas —con scores individuales entre 0.2606 y 0.2638— mediante el algoritmo de \emph{rank-based blending} descrito en la Sección~\ref{sec:metodologia} (Ecuación~\ref{eq:ensemble}), buscando aprovechar la diversidad entre modelos para extraer señal incremental.

% -----------------------------------------------------------
\subsubsection{Configuración del ensemble y scores individuales}

Se ensamblaron $K = 7$ modelos base cuyos scores individuales en el leaderboard público oscilan entre 0.2606 y 0.2638. La Tabla~\ref{tab:ensemble} detalla el score y el peso base asignado a cada uno.

\begin{table}[h!]
  \caption{Modelos base del ensemble. Los pesos $w_k$ suman 1.00. El modelo dominante concentra el 65\% del peso total.}
  \label{tab:ensemble}
  \centering
  \begin{tabular}{lcc}
    \toprule
    Modelo & Score público & Peso base $w_k$ \\
    \midrule
    \texttt{submission0.2638}        & 0.2638 & 0.65 \\
    \texttt{submissionhedgefund0.2637} & 0.2637 & 0.14 \\
    \texttt{submissionhedgefund0.2636} & 0.2636 & 0.08 \\
    \texttt{submissionhedgefund0.2635} & 0.2635 & 0.06 \\
    \texttt{submissionhedgefund0.2620} & 0.2620 & 0.04 \\
    \texttt{submissionhedgefund0.2608} & 0.2608 & 0.02 \\
    \texttt{submissionhedgefund0.2606} & 0.2606 & 0.01 \\
    \midrule
    \textbf{Ensemble} (Ecuación~\ref{eq:ensemble}) & \textbf{0.2639} & --- \\
    \bottomrule
  \end{tabular}
\end{table}

% -----------------------------------------------------------
\subsubsection{Resultado en el leaderboard}

El ensemble (Ecuación~\ref{eq:ensemble}) alcanzó un \emph{Weighted RMSE Skill Score} de \textbf{0.2639} en el leaderboard público, superando al mejor modelo individual (0.2638) en $\Delta = +0.0001$. Aunque la mejora absoluta es marginal, valida que la diversidad entre modelos, capturada por la mezcla de rangos bidireccional, aporta señal incremental incluso en un régimen de SNR extremadamente bajo.

% -----------------------------------------------------------
\subsubsection{Estadísticas descriptivas de las predicciones finales}

La Tabla~\ref{tab:stats} resume la distribución de las 1\,447\,107 predicciones generadas por el ensemble sobre el conjunto de prueba.

\begin{table}[h!]
  \caption{Estadísticas descriptivas del vector de predicción final del ensemble.}
  \label{tab:stats}
  \centering
  \begin{tabular}{lr}
    \toprule
    Estadístico & Valor \\
    \midrule
    Observaciones   & 1\,447\,107 \\
    Media           & $-0.886$ \\
    Desv. estándar  & $4.067$ \\
    Mínimo          & $-73.92$ \\
    Percentil 25    & $-0.137$ \\
    Mediana         & $-0.002$ \\
    Percentil 75    & $\phantom{-}0.000$ \\
    Máximo          & $\phantom{-}12.34$ \\
    \bottomrule
  \end{tabular}
\end{table}

La distribución muestra una marcada asimetría negativa (media $\ll$ mediana), coherente con la naturaleza de las series financieras anonimizadas: la mayoría de predicciones se concentran en torno a cero, con colas pesadas a la izquierda generadas por los horizontes de mayor duración ($h = 10, 25$).


% ============================================================
% REFERENCIAS
% ============================================================

\bibliographystyle{abbrvnat}
\bibliography{references}


\end{document}